\bookchapter{บทที่ 52: คอมโพเนนต์ทั้งสี่กับวงจรชีวิตของ Activity}

\chapterintro{ถ้า Compose คือ “ใบหน้า” ของแอป บทนี้ก็คือเรื่อง “โครงกระดูก” ของแอป ระบบ Android มีคอมโพเนนต์พื้นฐานสี่ตัว แอปทุกตัวประกอบขึ้นจากมัน และในบรรดานั้น Activity สำคัญที่สุด มี “แผนภาพวงจรชีวิต” ที่นักพัฒนา Android ทุกคนต้องสลักไว้ในหัว จบบทนี้ คุณจะเข้าใจปริศนาอมตะอย่าง “ทำไมสลับไปแอปอื่นแล้วกลับมา สถานะในแอปฉันหายหมด”}

\booksection{คอมโพเนนต์ทั้งสี่: “เสาหลัก” ของ Android}

ทางการ Android กำหนดคอมโพเนนต์พื้นฐานไว้สี่แบบ แบ่งกันรับผิดชอบงานคนละประเภท:

\begin{center}
\begin{tabularx}{\linewidth}{|X|X|X|}
\hline
\textbf{คอมโพเนนต์} & \textbf{หน้าที่} & \textbf{เปรียบเทียบง่าย ๆ} \tabularnewline
\hline
Activity & หน้าจอที่มองเห็นได้ & หน้าร้าน — ผู้ใช้มองเห็นและโต้ตอบได้ \tabularnewline
Service & "คนล่องหน" ทำงานเบื้องหลัง & คนงานคลัง — เปิดเพลง ดาวน์โหลดไฟล์ ไม่ต้องมีหน้าจอ \tabularnewline
BroadcastReceiver & รับการแจ้งจากระบบ (broadcast) & ห้องรับสาร — ลำโพงประกาศ "แบตหมด" "เน็ตหลุด" \tabularnewline
ContentProvider & แชร์ข้อมูลข้ามแอป & ตู้เอกสารส่วนกลาง — เช่น แชร์รายชื่อผู้ติดต่อให้แอปอื่น \tabularnewline
\hline
\end{tabularx}
\end{center}

\begin{center}\small\color{gray}ภาพรวมคอมโพเนนต์ทั้งสี่ของ Android\end{center}

ช่วงมือใหม่ คุณใช้เวลา 95\% กับ Activity ส่วน Service กับ BroadcastReceiver ใช้ตอนทำเล่นเพลง ส่งการแจ้งเตือน และ ContentProvider ใช้น้อยที่สุด (ระบบห่อข้อมูลที่ใช้บ่อยไว้ให้แล้ว) อีกเรื่องต้องระวัง: คอมโพเนนต์ทั้งสี่ต้อง “ขึ้นทะเบียน” ใน AndroidManifest.xml ด้วย ไม่ขึ้นทะเบียนระบบก็หาไม่เจอ

\booksection{วงจรชีวิตของ Activity: “แผนที่ขุมทรัพย์” แผนแรกของนักพัฒนา Android}

แอปบนมือถือถูกขัดจังหวะได้ทุกเมื่อ: สายเรียกเข้า ผู้ใช้สลับไปตอบแชท หน่วยความจำระบบไม่พอเลยโดนฆ่าทิ้ง... ชีวิตหนึ่งของ Activity ก็คือการเดินทางวนไปมาบนแผนภาพนี้:

onCreate (เกิด: เริ่มต้น)

onStart (มองเห็นแล้ว)

onResume (อยู่พื้นหน้า กำลังโต้ตอบ)

onPause (ถูกบังบางส่วน)

onStop (มองไม่เห็นเลย)

onDestroy (ทำลาย)

onRestart / onStart / onResume: กลับพื้นหน้าด้วยเส้นทางเดิม

\begin{center}\small\color{gray}ภาพที่ 4: วงจรชีวิตของ Activity (ไปข้างหน้า: แถวบนเกิด แถวล่างออกจากฉาก ตอนกลับก็ย้อนเส้นทางเดิม)\end{center}

\begin{center}
\begin{tabularx}{\linewidth}{|X|X|X|}
\hline
\textbf{Callback} & \textbf{เกิดขึ้นเมื่อไร} & \textbf{การใช้งานทั่วไป} \tabularnewline
\hline
onCreate & ตอน Activity ถูกสร้าง (ครั้งเดียว) & โหลดหน้าจอ กำหนดค่าเริ่มต้นข้อมูล \tabularnewline
onStart & หน้าจอกลายเป็นมองเห็นได้ & เริ่มตรรกะเบา ๆ บางอย่าง \tabularnewline
onResume & ได้โฟกัส พร้อมโต้ตอบ & เริ่มอนิเมชัน เล่นต่อ \tabularnewline
onPause & ถูกบังบางส่วน/เพิ่งถูกสลับออก & หยุดอนิเมชันชั่วคราว บันทึกฉบับร่าง \tabularnewline
onStop & มองไม่เห็นเลยสักนิด & คืนทรัพยากรหนัก พักการดาวน์โหลด \tabularnewline
onDestroy & ก่อนถูกทำลาย & การเก็บกวาดครั้งสุดท้าย \tabularnewline
\hline
\end{tabularx}
\end{center}

\begin{center}\small\color{gray}การใช้งานทั่วไปของ callback วงจรชีวิตทั้งหก\end{center}

\textbf{ข้อควรระวัง: หมุนหน้าจอจะ “ฆ่า” Activity}\par

ฉากพลาดคลาสสิกที่สุดของมือใหม่: หมุนหน้าจอทีเดียว ข้อความในช่องกรอกหายหมด เหตุผล: หมุนหน้าจอโดยค่าเริ่มต้นจะทำลายและสร้าง Activity ใหม่ (เดินผ่าน onDestroy onCreate หนึ่งรอบ) ทางแก้: ใช้ rememberSaveable เก็บสถานะ (ตัวแทนของ remember ใน Compose) มันจะกู้ค่ากลับมาให้เองตอนสร้างใหม่ จะได้ใช้ซ้ำอยู่บ่อย ๆ ในบทที่ 6 กับ 7

\booksection{Intent: “ตั๋วโดยสาร” เดินทางเข้าออกแอป}

Intent (เจตนา) คือ “ตั๋ว” ที่ใช้สื่อสารระหว่างคอมโพเนนต์ อยากกระโดดจากหน้าหนึ่งไปอีกหน้า หรือเปิดเบราว์เซอร์ เรียกกล้อง ก็ต้อง “ซื้อตั๋วขึ้นรถ” ทั้งนั้น:

\begin{lstlisting}[language=Kotlin]
// แบบที่หนึ่ง: Intent แบบชัด—ระบุปลายทางเองเลย
val intent = Intent(this, DetailActivity::class.java)
intent.putExtra("name", "ซานผวน")     // เหมือนฝาก "สัมภาระ" ขึ้นรถ (ข้อมูล)
startActivity(intent)

// แบบที่สอง: Intent แบบนัย—บอกแค่จะทำอะไร ระบบหาตัวทำเอง
val web = Intent(Intent.ACTION_VIEW, Uri.parse("https://developer.android.com"))
startActivity(web)   // ระบบจะเปิดเบราว์เซอร์/หรือให้คุณเลือก
\end{lstlisting}

สองแบบมีหน้าที่คนละแบบ: แบบชัดเจน (explicit) ใช้กระโดดหน้าภายในแอปตัวเอง ส่วนแบบโดยนัย (implicit) ใช้เรียกความสามารถของระบบหรือแอปอื่น (เปิดเว็บ โทรออก แชร์) ในโปรเจกต์ Compose อยู่ใน Composable ใช้ LocalContext.current ดึง Context มาแล้วเรียก startActivity ได้เลย

\booksection{หลายหน้ากับการนำทาง}

แอปจริง ๆ มีหลายหน้า วิธีนำทางที่ทางการแนะนำสำหรับ Compose คือ Navigation Compose (ไลบรารี androidx.navigation):

\begin{lstlisting}[language=Kotlin]
val navController = rememberNavController()

NavHost(navController, startDestination = "home") {
    composable("home") { HomeScreen(navController) }
    composable("detail/{name}") { entry ->
        val name = entry.arguments?.getString("name")
        DetailScreen(name ?: "ไม่ทราบ")
    }
}

// เปลี่ยนหน้าในหน้าจอ:
navController.navigate("detail/ซานผวน")     // พร้อมพารามิเตอร์
navController.popBackStack()              // กลับหน้าก่อนหน้า
\end{lstlisting}

\begin{tcbinfo}{เคล็ดลับ: หน้าเยอะขึ้นแล้วทำยังไงดี?}
พอโปรเจกต์โตขึ้น ให้รวมเส้นทาง (route) ของทุกหน้าไว้ในไฟล์เดียว (เช่น AppNavGraph.kt) จัดการแบบศูนย์กลาง อย่าปล่อยให้ตรรกะ “จะกระโดดไปไหน” กระจัดกระจายอยู่ตามแต่ละหน้า นี่คือก้าวแรกจาก “รันได้” สู่ “ดูแลต่อได้”

\end{tcbinfo}

\booksection{รู้จักไว้ระหว่างทาง: Fragment คืออะไร (และทำไมหนังสือเล่มนี้ไม่สอน)}

ค้นหาบทเรียนเก่าเมื่อไร คุณต้องเจอ Fragment แน่ ๆ — มันคือ “ชิ้นส่วนหน้าจอที่ใช้ซ้ำได้” อยู่ใน Activity เป็นพระเอกตัวจริงในสถาปัตยกรรมหลายหน้าแบบเก่า แต่ยุค Compose ทำหน้าด้วยฟังก์ชัน Composable + เส้นทาง Navigation ก็พอแล้ว ทางการก็ลดบทบาท Fragment ลง หนังสือเล่มนี้จึงไม่ใส่มันไว้ในเส้นทางหลัก

คำแนะนำสำหรับคุณ: โปรเจกต์ใหม่ไม่ต้องไปใช้ Fragment เอง แต่พออ่านโค้ดเก่าหรือโดนถามตอนสัมภาษณ์ แค่รู้ว่ามันคือ “ชิ้นส่วนหน้าจอที่ฝังอยู่ใน Activity” ก็พอ แหล่งเรียนรู้ให้เลือกที่เป็นหลังปี 2023 เป็นอันดับแรกเสมอ — “ปฏิทินเก่า” ของสแตกเทคโนโลยี Android อัปเดตเร็วกว่าที่คิด

\begin{tcbpractice}{ลงมือทำด้วยกัน: ใส่ "ความจำ" ให้ตัวนับ}
ปรับตัวนับจากบทที่ 51: 1) เปลี่ยน remember เป็น rememberSaveable แล้วหมุนหน้าจอดูว่าเลขยังอยู่ไหม; 2) เพิ่มปุ่ม “เกี่ยวกับ” ใช้ Navigation กระโดดไปหน้าที่สองแสดงหมายเลขเวอร์ชัน; 3) พิมพ์ log (Log.d) ใน onCreate/onPause ของทั้งสองหน้า แล้วใช้ Logcat ดูลำดับการเรียกวงจรชีวิต ทำครบสามข้อนี้ วงจรชีวิตก็ไม่ใช่แค่ทฤษฎีบนกระดาษอีกต่อไป

\end{tcbpractice}

\begin{tcbsummary}{สรุปบทนี้}
\begin{itemize}
  \item คอมโพเนนต์ทั้งสี่: Activity (หน้าร้าน), Service (เบื้องหลัง), BroadcastReceiver (ลำโพง), ContentProvider (ตู้เอกสาร)
  \item callback วงจรชีวิตหกตัว: onCreate onStart onResume onPause onStop onDestroy
  \item หมุนหน้าจอจะทำลายและสร้าง Activity ใหม่: ใช้ rememberSaveable ช่วยรักษาสถานะ
  \item Intent คือ "ตั๋ว": แบบชัดเจนกระโดดหน้าในแอปตัวเอง แบบโดยนัยเรียกใช้ความสามารถของระบบ
  \item หลายหน้าใช้ Navigation Compose: เส้นทาง NavHost + composable
\end{itemize}

\end{tcbsummary}
